\documentclass[11pt]{article}
\usepackage[margin=1in]{geometry}
\usepackage{booktabs}
\usepackage{longtable}
\usepackage{hyperref}
\usepackage{array}
\title{Behavior, Uncertainty, and Meta-Control\\A Working Note}
\author{Seungpil Lee}
\date{\today}

\begin{document}
\maketitle

\section{Introduction}
This working note tracks a small analysis program that connects two external perspectives to the
current Meta-CoT project. The first perspective is a four-behavior view of self-improving
reasoners: verification, backtracking, subgoal setting, and backward chaining. The second
perspective is an uncertainty view in which useful reasoning depends on epistemic verbalization
rather than procedural text alone.

\section{Why This Note Exists}
The main training track now uses three nodes for control-v5 SFT and later RL. In parallel, this note
defines a separate analysis track that can interpret model outputs before and after training. The main
question is whether current meta traces already behave like control actions conditioned on uncertainty.

\section{Current Analysis Scope}
The first pass uses locally available evaluation outputs for existing Meta-CoT models. The placeholder
goal is to measure benchmark-level rates for verification, redirect-like behavior, subgoal language,
backward-chaining language, uncertainty language, diagnosis language, and epistemic behavior.

\section{Planned Outputs}
\begin{itemize}
\item A summary table over available models and benchmarks.
\item A small qualitative slice on difficult AIME examples.
\item A decision memo for how to map the analysis back into reward design and curriculum/RAG.
\end{itemize}

\section{Status}
This PDF is a placeholder shell. The quantitative table and qualitative examples are generated in the
paired Markdown report and can be folded into a later revision once more model outputs are available.

\end{document}
